% --- IMPORTS ---
\documentclass[leqno]{article}
\usepackage{graphicx}
\usepackage{dsfont}
\usepackage{mathtools}
\usepackage{amssymb}
\usepackage{amsmath}
\usepackage{amsthm}
\usepackage{float}
\usepackage{ragged2e}
\usepackage{tensor}
\usepackage{fancyhdr}
\usepackage{empheq}
\usepackage[most]{tcolorbox}
\usepackage{enumitem}
\usepackage{dirtytalk}
\usepackage{marginnote}
\usepackage{microtype}
\usepackage[
  a4paper,
  left=0.8in,
  right=1in,
  top=1in,
  bottom=0.5in,
  headheight=14pt,
  headsep=0.4in
]{geometry}
\setlength{\parindent}{0cm}
% --- TikZ Packages ---
\usepackage{pgfplots}
\pgfplotsset{compat=1.18}
\usepgfplotslibrary{fillbetween, polar}
\usetikzlibrary{calc, positioning}
\usetikzlibrary{angles, quotes}
\usetikzlibrary{arrows.meta}
\usetikzlibrary{decorations.pathreplacing, decorations.markings}
\usetikzlibrary{patterns}
\usetikzlibrary{intersections}
% --- URL Packages ---
\usepackage[colorlinks=true,linkcolor=red,citecolor=magenta,urlcolor=black]{hyperref}%
\urlstyle{same}
% --- END OF IMPORTS ---

% --- CUSTOM COMMANDS ---
\RenewDocumentCommand{\marks}{o m}{%
  \IfValueTF{#1}%
    {\marginnote{\raggedright\textbf{[#2]}}[#1]}%
    {\marginnote{\raggedright\textbf{[#2]}}}%
}
\newcommand{\vspacerule}[2]{%
  \par\vspace*{#1}%
  \noindent\hrulefill\par
  \vspace*{#2}%
}
\definecolor{scarlet}{RGB}{205, 40, 75}
\newcommand{\questionitem}{%
  \item\phantomsection
  \addcontentsline{toc}{section}{Question \number\value{enumi}}%
}
\renewcommand{\contentsname}{Questions}
% --- END OF CUSTOM COMMANDS ---

% --- HEADER CONFIG ---
\pagestyle{fancy}
\fancyhead{}
\fancyfoot{}
\renewcommand{\headrulewidth}{0.9pt}
\lhead{\textit{Solutions} - 2nd Order Differential Equations}
\chead{}
\rhead{\href{https://danieljsmith.org}{danieljsmith.org}}
% --- END OF HEADER CONFIG ---

\begin{document}
\vspace*{20pt}
\tableofcontents
\newpage
\begin{enumerate}[
  label=\textbf{\arabic*.},
  leftmargin=*,
  topsep=0pt,
  partopsep=0pt,
  parsep=0pt
]

% ---- Question 1 ----
\questionitem
% [Source: _QBT___Solns__2nd_Order_Differential_Equations, Q1]

\begin{enumerate}[label=\textbf{(\alph*)}]
  \item Show that \[3\cosh x - \sinh x = e^x + 2e^{-x}\] \marks{2} \\
  \item Hence solve the differential equation
  \[
  \frac{\mathrm{d}^2 y}{\mathrm{d}x^2} - \frac{\mathrm{d}y}{\mathrm{d}x} - 2y = 6(3\cosh x - \sinh x)
  \]
  given that, when $x=0$, $y=1$ and $\frac{\mathrm{d}y}{\mathrm{d}x}=1$. \marks{10} \\
\end{enumerate}

\vspace*{10pt}

\begin{tcolorbox}[colback=white, colframe=scarlet, title={\textbf{Solution}}, breakable, grow to left by=6mm, grow to right by=10mm]
\begin{enumerate}[label=\textbf{(\alph*)}]
  \item Using
\[
\cosh x=\frac{e^x+e^{-x}}{2}
\qquad\text{and}\qquad
\sinh x=\frac{e^x-e^{-x}}{2}
\]
we have
\begin{align*}
3\cosh x-\sinh x
&=3\left(\frac{e^x+e^{-x}}{2}\right)-\left(\frac{e^x-e^{-x}}{2}\right)\\
&=\frac{3e^x+3e^{-x}-e^x+e^{-x}}{2}\\
&=\frac{2e^x+4e^{-x}}{2}\\
&=e^x+2e^{-x}
\end{align*}
Hence
\[
3\cosh x-\sinh x=e^x+2e^{-x}
\]

  \item From part (a),
\[
6(3\cosh x-\sinh x)=6(e^x+2e^{-x})=6e^x+12e^{-x}
\]
so the differential equation becomes
\[
\frac{\mathrm{d}^2 y}{\mathrm{d}x^2}-\frac{\mathrm{d}y}{\mathrm{d}x}-2y=6e^x+12e^{-x}
\]

First solve the homogeneous equation
\[
\frac{\mathrm{d}^2 y}{\mathrm{d}x^2}-\frac{\mathrm{d}y}{\mathrm{d}x}-2y=0
\]
The auxilliary equation is
\[
\lambda^2-\lambda-2=0
\]
Factorising,
\begin{align*}
\lambda^2-\lambda-2&=0\\
(\lambda-2)(\lambda+1)&=0
\end{align*}
So the roots are $\lambda=2$ and $\lambda=-1$. Therefore the complementary function is
\[
y_c=Ae^{2x}+Be^{-x}
\]

Now find a particular integral. Since the right-hand side is $6e^x+12e^{-x}$, take
\[
y_p=Ce^x+Dxe^{-x}
\]
We use $Dxe^{-x}$ rather than $De^{-x}$ because $e^{-x}$ is already part of the complementary function.

For the $Ce^x$ term,
\[
\frac{\mathrm{d}}{\mathrm{d}x}(Ce^x)=Ce^x,
\qquad
\frac{\mathrm{d}^2}{\mathrm{d}x^2}(Ce^x)=Ce^x
\]
Substituting into the left-hand side,
\begin{align*}
\frac{\mathrm{d}^2}{\mathrm{d}x^2}(Ce^x)-\frac{\mathrm{d}}{\mathrm{d}x}(Ce^x)-2(Ce^x)
&=Ce^x-Ce^x-2Ce^x\\
&=-2Ce^x
\end{align*}
This must equal $6e^x$, so
\[
-2C=6
\]
Hence
\[
C=-3
\]

For the $Dxe^{-x}$ term,
\begin{align*}
y&=Dxe^{-x}\\
\frac{\mathrm{d}y}{\mathrm{d}x}
&=D\frac{\mathrm{d}}{\mathrm{d}x}(xe^{-x})\\
&=D\left(e^{-x}-xe^{-x}\right)\\
&=D(1-x)e^{-x}
\end{align*}
Differentiate again:
\begin{align*}
\frac{\mathrm{d}^2 y}{\mathrm{d}x^2}
&=D\frac{\mathrm{d}}{\mathrm{d}x}\left((1-x)e^{-x}\right)\\
&=D\left((-1)e^{-x}+(1-x)(-e^{-x})\right)\\
&=D(x-2)e^{-x}
\end{align*}
Substituting into the left-hand side,
\begin{align*}
\frac{\mathrm{d}^2 y}{\mathrm{d}x^2}-\frac{\mathrm{d}y}{\mathrm{d}x}-2y
&=D(x-2)e^{-x}-D(1-x)e^{-x}-2Dxe^{-x}\\
&=D\bigl(x-2-1+x-2x\bigr)e^{-x}\\
&=-3De^{-x}
\end{align*}
This must equal $12e^{-x}$, so
\[
-3D=12
\]
Hence
\[
D=-4
\]

Therefore a particular integral is
\[
y_p=-3e^x-4xe^{-x}
\]
and the general solution is
\[
y=Ae^{2x}+Be^{-x}-3e^x-4xe^{-x}
\]

Now use the initial conditions.

When $x=0$, $y=1$:
\begin{align*}
1
&=Ae^0+Be^0-3e^0-4(0)e^0\\
&=A+B-3
\end{align*}
So
\[
A+B=4
\]

Differentiate the general solution:
\begin{align*}
\frac{\mathrm{d}y}{\mathrm{d}x}
&=2Ae^{2x}-Be^{-x}-3e^x-4\frac{\mathrm{d}}{\mathrm{d}x}(xe^{-x})\\
&=2Ae^{2x}-Be^{-x}-3e^x-4(1-x)e^{-x}
\end{align*}
When $x=0$, $\frac{\mathrm{d}y}{\mathrm{d}x}=1$:
\begin{align*}
1
&=2Ae^0-Be^0-3e^0-4(1-0)e^0\\
&=2A-B-3-4\\
&=2A-B-7
\end{align*}
So
\[
2A-B=8
\]

Solve the simultaneous equations
\begin{align*}
A+B&=4\\
2A-B&=8
\end{align*}
Adding gives
\[
3A=12
\]
so
\[
A=4
\]
Then
\[
B=4-A=0
\]

Substituting back,
\[
y=4e^{2x}-3e^x-4xe^{-x}
\]

Therefore the required solution is
\[
y=4e^{2x}-3e^x-4xe^{-x}
\]
\end{enumerate}
\end{tcolorbox}


\newpage

% ---- Question 2 ----
\questionitem
% [Source: _QBT___Solns__2nd_Order_Differential_Equations, Q2]

Find the particular solution of the differential equation \\
\[
\frac{\mathrm{d}^2 y}{\mathrm{d}t^2}+4\frac{\mathrm{d}y}{\mathrm{d}t}+13y=8e^{-2t}\sin t
\] \\
given that, when $t=0$, $y=1$ and $\frac{\mathrm{d}y}{\mathrm{d}t}=2$. \marks{10}

\vspace*{10pt}

\begin{tcolorbox}[colback=white, colframe=scarlet, title={\textbf{Solution}}, breakable, grow to left by=6mm, grow to right by=10mm]
For the homogeneous part,
\[
\frac{\mathrm{d}^2 y}{\mathrm{d}t^2}+4\frac{\mathrm{d}y}{\mathrm{d}t}+13y=0
\]

The auxiliary equation is
\[
\lambda^2+4\lambda+13=0
\]

So
\[
\lambda=\frac{-4\pm\sqrt{16-52}}{2}=-2\pm 3i
\]

Hence
\[
y_c=e^{-2t}\left(A\cos 3t+B\sin 3t\right)
\]

For a particular integral, try
\[
y_p=e^{-2t}\left(a\sin t+b\cos t\right)
\]

Then
\begin{align*}
\frac{\mathrm{d}y_p}{\mathrm{d}t}
&= \frac{\mathrm{d}}{\mathrm{d}t}\left[e^{-2t}\left(a\sin t+b\cos t\right)\right] \\
&= -2e^{-2t}\left(a\sin t+b\cos t\right)
+e^{-2t}\left(a\cos t-b\sin t\right) \\
&= e^{-2t}\left[(-b-2a)\sin t+(a-2b)\cos t\right]
\end{align*}

Also,
\begin{align*}
\frac{\mathrm{d}^2y_p}{\mathrm{d}t^2}
&= \frac{\mathrm{d}}{\mathrm{d}t}\left\{e^{-2t}\left[(-b-2a)\sin t+(a-2b)\cos t\right]\right\} \\
&= -2e^{-2t}\left[(-b-2a)\sin t+(a-2b)\cos t\right] \\
&\qquad +e^{-2t}\left[(-b-2a)\cos t-(a-2b)\sin t\right] \\
&= e^{-2t}\left[(3a+4b)\sin t+(-4a+3b)\cos t\right]
\end{align*}

Substituting into
\[
\frac{\mathrm{d}^2 y}{\mathrm{d}t^2}+4\frac{\mathrm{d}y}{\mathrm{d}t}+13y=8e^{-2t}\sin t
\]

gives
\begin{align*}
& e^{-2t}\left[(3a+4b)\sin t+(-4a+3b)\cos t\right] \\
& \qquad +4e^{-2t}\left[(-b-2a)\sin t+(a-2b)\cos t\right] \\
& \qquad +13e^{-2t}\left(a\sin t+b\cos t\right) \\
&= 8e^{-2t}\sin t
\end{align*}

So
\[
e^{-2t}\left(8a\sin t+8b\cos t\right)=8e^{-2t}\sin t
\]

Hence
\[
8a\sin t+8b\cos t=8\sin t
\]

Comparing coefficients,
\[
8a=8, \qquad 8b=0
\]

So
\[
a=1, \qquad b=0
\]

Therefore
\[
y_p=e^{-2t}\sin t
\]

Thus the general solution is
\[
y=e^{-2t}\left(A\cos 3t+B\sin 3t\right)+e^{-2t}\sin t
\]

So
\[
y=e^{-2t}\left(A\cos 3t+B\sin 3t+\sin t\right)
\]

Now use the initial conditions

Given $y(0)=1$,
\[
1=e^{0}\left(A\cos 0+B\sin 0+\sin 0\right)
\]

So
\[
1=A
\]

Hence
\[
A=1
\]

Now differentiate
\begin{align*}
\frac{\mathrm{d}y}{\mathrm{d}t}
&= \frac{\mathrm{d}}{\mathrm{d}t}\left[e^{-2t}\left(A\cos 3t+B\sin 3t+\sin t\right)\right] \\
&= -2e^{-2t}\left(A\cos 3t+B\sin 3t+\sin t\right) \\
&\qquad +e^{-2t}\left(-3A\sin 3t+3B\cos 3t+\cos t\right) \\
&= e^{-2t}\left(-3A\sin 3t+3B\cos 3t+\cos t-2A\cos 3t-2B\sin 3t-2\sin t\right)
\end{align*}

At $t=0$,
\[
\frac{\mathrm{d}y}{\mathrm{d}t}(0)=3B+1-2A
\]

Since $\frac{\mathrm{d}y}{\mathrm{d}t}(0)=2$ and $A=1$,
\[
2=3B+1-2
\]

So
\[
2=3B-1
\]

Hence
\[
3B=3
\]

So
\[
B=1
\]

Substituting $A=1$ and $B=1$ into the solution,
\[
y=e^{-2t}\bigl(\cos 3t+\sin 3t+\sin t\bigr)
\]

Hence the required particular solution is
\[
y=e^{-2t}\bigl(\cos 3t+\sin 3t+\sin t\bigr)
\]
\end{tcolorbox}


\newpage

% ---- Question 3 ----
\questionitem
% [Source: _QBT___Solns__2nd_Order_Differential_Equations, Q3]

A damped oscillator is governed by \\
\[
\frac{\mathrm{d}^2 y}{\mathrm{d}t^2}+6\frac{\mathrm{d}y}{\mathrm{d}t}+13y=75\sin 2t
\] \\
\begin{enumerate}[label=\textbf{(\alph*)}]
  \item Find the general solution of this differential equation. \marks{6} \\
  \item Show that, whatever the initial conditions, the motion is approximately of the form
  \[
  y\approx R\sin(2t-\phi)
  \]
  for large positive values of $t$, where $R>0$ and $0<\phi<\frac{\pi}{2}$. Determine $R$ and $\phi$. \marks{4} \\
\end{enumerate}

\vspace*{10pt}

\begin{tcolorbox}[colback=white, colframe=scarlet, title={\textbf{Solution}}, breakable, grow to left by=6mm, grow to right by=10mm]
\begin{enumerate}[label=\textbf{(\alph*)}]
  \item First solve the associated homogeneous equation
  \[
  \frac{\mathrm{d}^2y}{\mathrm{d}t^2}+6\frac{\mathrm{d}y}{\mathrm{d}t}+13y=0
  \]

  Using the trial solution $y=\mathrm{e}^{\lambda t}$ gives the auxiliary equation
  \[
  \lambda^2+6\lambda+13=0
  \]

  Solving,
  \begin{align*}
  \lambda&=\frac{-6\pm\sqrt{6^2-4(1)(13)}}{2} \\[4pt]
   &=-3\pm 2\mathrm{i}
  \end{align*}

  Hence the complementary function is
  \[
  y_c=\mathrm{e}^{-3t}(A\cos 2t+B\sin 2t)
  \]

  For a particular integral, try
  \[
  y_p=a\sin 2t+b\cos 2t
  \]

  Then
  \begin{align*}
  \frac{\mathrm{d}y_p}{\mathrm{d}t}&=2a\cos 2t-2b\sin 2t \\
  \frac{\mathrm{d}^2y_p}{\mathrm{d}t^2}&=-4a\sin 2t-4b\cos 2t
  \end{align*}

  Substituting into the differential equation:
  \begin{align*}
  &\left(-4a\sin 2t-4b\cos 2t\right)+6\left(2a\cos 2t-2b\sin 2t\right)+13\left(a\sin 2t+b\cos 2t\right) \\
  &=75\sin 2t
  \end{align*}

  Collecting $\sin 2t$ and $\cos 2t$ terms,
  \begin{align*}
  &( -4a-12b+13a)\sin 2t+(12a-4b+13b)\cos 2t \\
  &=75\sin 2t
  \end{align*}

  So
  \[
  (9a-12b)\sin 2t+(12a+9b)\cos 2t=75\sin 2t
  \]

  Equating coefficients gives
  \begin{align*}
  9a-12b&=75 \\
  12a+9b&=0
  \end{align*}

  Solve these simultaneously. From the second equation,
  \[
  4a+3b=0 \quad \Rightarrow \quad b=-\frac{4a}{3}
  \]

  Substitute into the first:
  \begin{align*}
  9a-12\left(-\frac{4a}{3}\right)&=75 \\
  9a+16a&=75 \\
  25a&=75 \\
  a&=3
  \end{align*}

  Then
  \[
  b=-\frac{4(3)}{3}=-4
  \]

  Thus
  \[
  y_p=3\sin 2t-4\cos 2t
  \]

  Therefore the general solution is
  \[
  y=\mathrm{e}^{-3t}(A\cos 2t+B\sin 2t)+3\sin 2t-4\cos 2t
  \]

  \item From part (a),
  \[
  y=\mathrm{e}^{-3t}(A\cos 2t+B\sin 2t)+3\sin 2t-4\cos 2t
  \]

  As $t\to+\infty$, the exponential factor tends to $0$. Also $\sin 2t$ and $\cos 2t$ remain bounded, so
  \begin{align*}
  \mathrm{e}^{-3t}(A\cos 2t+B\sin 2t) 
  &\longrightarrow 0
  \end{align*}

  Therefore, whatever the initial conditions,
  \[
  y\approx 3\sin 2t-4\cos 2t
  \]
  for large positive $t$.

  Now write this in the form $R\sin(2t-\phi)$. Using the identity
  \[
  \sin(2t-\phi)=\sin 2t\cos\phi-\cos 2t\sin\phi
  \]

  we have
  \[
  R\sin(2t-\phi)=R\cos\phi\,\sin 2t-R\sin\phi\,\cos 2t
  \]

  Comparing with $3\sin 2t-4\cos 2t$ gives
  \begin{align*}
  R\cos\phi&=3 \\
  R\sin\phi&=4
  \end{align*}

  Squaring and adding,
  \begin{align*}
  R^2\cos^2\phi+R^2\sin^2\phi&=3^2+4^2 \\
  R^2(\cos^2\phi+\sin^2\phi)&=25 \\
  R^2&=25
  \end{align*}

  Since $R>0$,
  \[
  R=5
  \]

  Also,
  \[
  \tan\phi=\frac{R\sin\phi}{R\cos\phi}=\frac{4}{3}
  \]

  Since $0<\phi<\frac{\pi}{2}$,
  \[
  \phi=\arctan\left(\frac{4}{3}\right)\approx 0.927
  \]

  Hence for large positive $t$,
  \[
  y\approx 5\sin\left(2t-\arctan\frac{4}{3}\right)
  \]

  So the required constants are $R=5$ and $\phi=\arctan\left(\frac{4}{3}\right)\approx 0.927$
\end{enumerate}
\end{tcolorbox}


\newpage

% ---- Question 4 ----
\questionitem
% [Source: _QBT___Solns__2nd_Order_Differential_Equations, Q4]

A function $y(t)$ satisfies the differential equation \\
\[
\frac{\mathrm{d}^2 y}{\mathrm{d}t^2}+2\frac{\mathrm{d}y}{\mathrm{d}t}+5y=10\cos t
\] \\
Given that, when $t=0$, $y=5$ and $\frac{\mathrm{d}y}{\mathrm{d}t}=2$, find the particular solution. \marks{10}

\vspace*{10pt}

\begin{tcolorbox}[colback=white, colframe=scarlet, title={\textbf{Solution}}, breakable, grow to left by=6mm, grow to right by=10mm]
First solve the associated homogeneous equation
\[
\frac{\mathrm{d}^2 y}{\mathrm{d} t^2}+2\frac{\mathrm{d} y}{\mathrm{d} t}+5y=0
\]

Its auxiliary equation is
\[
\lambda^2+2\lambda+5=0
\]

Solving,
\begin{align*}
\lambda &= \frac{-2\pm\sqrt{2^2-4(1)(5)}}{2} \\[4pt]
  &= -1\pm 2i
\end{align*}

So the complementary function is
\[
y_c=e^{-t}(C_1\cos 2t+C_2\sin 2t)
\]

Now find a particular integral for
\[
\frac{\mathrm{d}^2 y}{\mathrm{d} t^2}+2\frac{\mathrm{d} y}{\mathrm{d} t}+5y=10\cos t
\]

Since the forcing term is $10\cos t$, try
\[
y_p=A\cos t+B\sin t
\]

Then
\[
\frac{\mathrm{d} y_p}{\mathrm{d} t}=-A\sin t+B\cos t
\]
and
\[
\frac{\mathrm{d}^2 y_p}{\mathrm{d} t^2}=-A\cos t-B\sin t
\]

Substituting into the differential equation,
\begin{align*}
(-A\cos t-B\sin t)+2(-A\sin t+B\cos t)+5(A\cos t+B\sin t)
&=10\cos t
\end{align*}

Collecting $\cos t$ and $\sin t$ terms,
\begin{align*}
(-A+2B+5A)\cos t+(-B-2A+5B)\sin t &= 10\cos t \\
(4A+2B)\cos t+(4B-2A)\sin t &= 10\cos t
\end{align*}

Equating coefficients gives
\[
4A+2B=10,\qquad 4B-2A=0
\]

From $4B-2A=0$,
\[
A=2B
\]

Substitute into $4A+2B=10$:
\begin{align*}
4(2B)+2B &= 10 \\
8B+2B &= 10 \\
10B &= 10 \\
B &= 1
\end{align*}

Hence
\[
A=2
\]

So the particular integral is
\[
y_p=2\cos t+\sin t
\]

Therefore the general solution is
\[
y=e^{-t}(C_1\cos 2t+C_2\sin 2t)+2\cos t+\sin t
\]

Now use the initial conditions.

When $t=0$, $y=5$:
\begin{align*}
5 &= e^0(C_1\cos 0+C_2\sin 0)+2\cos 0+\sin 0 \\
  &= C_1+2
\end{align*}

So
\[
C_1=3
\]

Next differentiate the general solution:
\begin{align*}
\frac{\mathrm{d} y}{\mathrm{d} t}
&= \frac{\mathrm{d}}{\mathrm{d} t}\left[e^{-t}(C_1\cos 2t+C_2\sin 2t)\right]+\frac{\mathrm{d}}{\mathrm{d} t}(2\cos t+\sin t)
\end{align*}

Using the product rule on the first term,
\begin{align*}
\frac{\mathrm{d} y}{\mathrm{d} t}
&= -e^{-t}(C_1\cos 2t+C_2\sin 2t)
+e^{-t}(-2C_1\sin 2t+2C_2\cos 2t)
-2\sin t+\cos t \\
&= e^{-t}\left[(-C_1+2C_2)\cos 2t+(-2C_1-C_2)\sin 2t\right]-2\sin t+\cos t
\end{align*}

Now use $\frac{\mathrm{d} y}{\mathrm{d} t}=2$ when $t=0$:
\begin{align*}
2&=e^0\left[(-C_1+2C_2)\cos 0+(-2C_1-C_2)\sin 0\right]-2\sin 0+\cos 0 \\
&=(-C_1+2C_2)+1
\end{align*}

Substitute $C_1=3$:
\begin{align*}
2 &= -3+2C_2+1 \\
2 &= 2C_2-2 \\
4 &= 2C_2 \\
C_2 &= 2
\end{align*}

Hence the required particular solution satisfying the given conditions is
\[
y(t)=e^{-t}(3\cos 2t+2\sin 2t)+2\cos t+\sin t
\]
\end{tcolorbox}


\newpage

% ---- Question 5 ----
\questionitem
% [Source: _QBT___Solns__2nd_Order_Differential_Equations, Q5]

\begin{enumerate}[label=\textbf{(\alph*)}]
  \item Find the general solution of the differential equation
  \[
  \frac{\mathrm{d}^2 y}{\mathrm{d}x^2}+4\frac{\mathrm{d}y}{\mathrm{d}x}+13y=0
  \]
  \marks[-20pt]{2}
  \item Hence find the general solution of the differential equation
  \[
  \frac{\mathrm{d}^2 y}{\mathrm{d}x^2}+4\frac{\mathrm{d}y}{\mathrm{d}x}+13y=e^{-2x}(6\cos x-8\sin x)
  \]
  \marks[-20pt]{4}
\end{enumerate}

\vspace*{10pt}

\begin{tcolorbox}[colback=white, colframe=scarlet, title={\textbf{Solution}}, breakable, grow to left by=6mm, grow to right by=10mm]
\begin{enumerate}[label=\textbf{(\alph*)}]
  \item For the homogeneous equation, use the auxiliary equation
\[
\lambda^2+4\lambda+13=0
\]
Using the quadratic formula,
\begin{align*}
\lambda&=\frac{-4\pm\sqrt{4^2-4(1)(13)}}{2} \\[4pt]
&=-2\pm 3i
\end{align*}
Hence the complementary solution is
\[
y=e^{-2x}\left(A\cos 3x+B\sin 3x\right)
\]
So the general solution is $y=e^{-2x}\left(A\cos 3x+B\sin 3x\right)$ \\

  \item From part (a), the complementary function is
\[
y_c=e^{-2x}\left(C\cos 3x+D\sin 3x\right)
\]

For a particular integral, try
\[
y_p=e^{-2x}\left(a\cos x+b\sin x\right)
\]

Then
\begin{align*}
\frac{\mathrm{d}y_p}{\mathrm{d}x}
&= \frac{\mathrm{d}}{\mathrm{d}x}\left[e^{-2x}\left(a\cos x+b\sin x\right)\right] \\
&= -2e^{-2x}\left(a\cos x+b\sin x\right)
+e^{-2x}\left(-a\sin x+b\cos x\right) \\
&= e^{-2x}\left[(b-2a)\cos x+(-a-2b)\sin x\right]
\end{align*}

Also,
\begin{align*}
\frac{\mathrm{d}^2y_p}{\mathrm{d}x^2}
&= \frac{\mathrm{d}}{\mathrm{d}x}\left\{e^{-2x}\left[(b-2a)\cos x+(-a-2b)\sin x\right]\right\} \\
&= -2e^{-2x}\left[(b-2a)\cos x+(-a-2b)\sin x\right] \\
&\qquad +e^{-2x}\left[(-b+2a)\sin x+(-a-2b)\cos x\right] \\
&= e^{-2x}\left[(3a-4b)\cos x+(4a+3b)\sin x\right]
\end{align*}

Substituting into
\[
\frac{\mathrm{d}^2 y}{\mathrm{d}x^2}+4\frac{\mathrm{d}y}{\mathrm{d}x}+13y=e^{-2x}(6\cos x-8\sin x)
\]

gives
\begin{align*}
& e^{-2x}\left[(3a-4b)\cos x+(4a+3b)\sin x\right] \\
& \qquad +4e^{-2x}\left[(b-2a)\cos x+(-a-2b)\sin x\right] \\
& \qquad +13e^{-2x}\left(a\cos x+b\sin x\right) \\
&= e^{-2x}(6\cos x-8\sin x)
\end{align*}

So
\[
e^{-2x}\left(8a\cos x+8b\sin x\right)=e^{-2x}(6\cos x-8\sin x)
\]

Hence
\[
8a\cos x+8b\sin x=6\cos x-8\sin x
\]

Comparing coefficients,
\[
8a=6, \qquad 8b=-8
\]

So
\[
a=\frac{3}{4}, \qquad b=-1
\]

Therefore
\[
y_p=e^{-2x}\left(\frac{3}{4}\cos x-\sin x\right)
\]

Thus the general solution is
\[
y=e^{-2x}\left(C\cos 3x+D\sin 3x\right)+e^{-2x}\left(\frac{3}{4}\cos x-\sin x\right)
\]

So
\[
y=e^{-2x}\left(C\cos 3x+D\sin 3x+\frac{3}{4}\cos x-\sin x\right)
\]
\end{enumerate}
\end{tcolorbox}


\newpage

% ---- Question 6 ----
\questionitem
% [Source: _QBT___Solns__2nd_Order_Differential_Equations, Q6]

\begin{enumerate}[label=\textbf{(\alph*)}]
  \item Show that \[\frac{\cosh x - \sinh x}{\cosh x + \sinh x}=e^{-2x}\] \marks{2} \\
  \item Hence find the particular solution of the differential equation
  \[
  \frac{\mathrm{d}^2 y}{\mathrm{d}x^2} + 4\frac{\mathrm{d}y}{\mathrm{d}x} + 4y = 4\left(\frac{\cosh x - \sinh x}{\cosh x + \sinh x}\right)
  \]
  given that, when $x=0$, $y=1$ and $\frac{\mathrm{d}y}{\mathrm{d}x}=1$. \marks{10} \\
\end{enumerate}

\vspace*{10pt}

\begin{tcolorbox}[colback=white, colframe=scarlet, title={\textbf{Solution}}, breakable, grow to left by=6mm, grow to right by=10mm]
\begin{enumerate}[label=\textbf{(\alph*)}]
  \item Using
\[
\cosh x=\frac{e^x+e^{-x}}{2}
\qquad \text{and} \qquad
\sinh x=\frac{e^x-e^{-x}}{2}
\]
we get
\begin{align*}
\cosh x-\sinh x
&= \frac{e^x+e^{-x}}{2}-\frac{e^x-e^{-x}}{2} \\
&= \frac{e^x+e^{-x}-e^x+e^{-x}}{2} \\
&= \frac{2e^{-x}}{2} \\
&= e^{-x}
\end{align*}
Also,
\begin{align*}
\cosh x+\sinh x
&= \frac{e^x+e^{-x}}{2}+\frac{e^x-e^{-x}}{2} \\
&= \frac{e^x+e^{-x}+e^x-e^{-x}}{2} \\
&= \frac{2e^x}{2} \\
&= e^x
\end{align*}
Therefore,
\[
\frac{\cosh x-\sinh x}{\cosh x+\sinh x}
= \frac{e^{-x}}{e^x}
= e^{-2x}
\]
Hence,
\[
\frac{\cosh x-\sinh x}{\cosh x+\sinh x}=e^{-2x}
\]

\item From part (a), the differential equation becomes
\[
\frac{\mathrm{d}^2 y}{\mathrm{d}x^2}+4\frac{\mathrm{d}y}{\mathrm{d}x}+4y=4e^{-2x}
\]

For the homogeneous equation,
\[
\frac{\mathrm{d}^2 y}{\mathrm{d}x^2}+4\frac{\mathrm{d}y}{\mathrm{d}x}+4y=0
\]

The auxiliary equation is
\[
\lambda^2+4\lambda+4=0
\]

So
\[
(\lambda+2)^2=0
\]

Hence the complementary function is
\[
y_c=(Cx+D)e^{-2x}
\]

Since the right-hand side is $4e^{-2x}$ and $-2$ is a repeated root of the auxiliary equation, try
\[
y_p=Ax^2e^{-2x}
\]

Differentiate:
\begin{align*}
\frac{\mathrm{d}y_p}{\mathrm{d}x}
&= \frac{\mathrm{d}}{\mathrm{d}x}\left(Ax^2e^{-2x}\right) \\
&= 2Axe^{-2x}-2Ax^2e^{-2x} \\
&= e^{-2x}(2Ax-2Ax^2)
\end{align*}

Differentiate again:
\begin{align*}
\frac{\mathrm{d}^2y_p}{\mathrm{d}x^2}
&= \frac{\mathrm{d}}{\mathrm{d}x}\left[e^{-2x}(2Ax-2Ax^2)\right] \\
&= -2e^{-2x}(2Ax-2Ax^2)+e^{-2x}(2A-4Ax) \\
&= e^{-2x}(2A-8Ax+4Ax^2)
\end{align*}

Substituting into the differential equation gives
\begin{align*}
\frac{\mathrm{d}^2y_p}{\mathrm{d}x^2}+4\frac{\mathrm{d}y_p}{\mathrm{d}x}+4y_p
&= e^{-2x}(2A-8Ax+4Ax^2)+4e^{-2x}(2Ax-2Ax^2)+4Ax^2e^{-2x} \\
&= e^{-2x}(2A)
\end{align*}

So
\[
2Ae^{-2x}=4e^{-2x}
\]

Hence
\[
2A=4
\]

So
\[
A=2
\]

Therefore
\[
y_p=2x^2e^{-2x}
\]

Thus the general solution is
\[
y=(Cx+D)e^{-2x}+2x^2e^{-2x}
\]

So
\[
y=e^{-2x}(2x^2+Cx+D)
\]

Now use the initial conditions

Since $y(0)=1$,
\begin{align*}
1
&= e^{0}(2\cdot 0^2+C\cdot 0+D) \\
&= D
\end{align*}

So
\[
D=1
\]

Now differentiate
\begin{align*}
\frac{\mathrm{d}y}{\mathrm{d}x}
&= \frac{\mathrm{d}}{\mathrm{d}x}\left[e^{-2x}(2x^2+Cx+D)\right] \\
&= -2e^{-2x}(2x^2+Cx+D)+e^{-2x}(4x+C) \\
&= e^{-2x}(4x+C-4x^2-2Cx-2D)
\end{align*}

At $x=0$,
\[
\frac{\mathrm{d}y}{\mathrm{d}x}(0)=C-2D
\]

Since $\frac{\mathrm{d}y}{\mathrm{d}x}(0)=1$ and $D=1$,
\[
1=C-2
\]

Hence
\[
C=3
\]

Substituting $C=3$ and $D=1$ into the solution gives
\[
y=(2x^2+3x+1)e^{-2x}
\]

Therefore the required particular solution is
\[
y=(2x^2+3x+1)e^{-2x}
\]
\end{enumerate}
\end{tcolorbox}


\newpage

% ---- Question 7 ----
\questionitem
% [Source: _QBT___Solns__2nd_Order_Differential_Equations, Q7]

Solve the differential equation \\
\[
\frac{\mathrm{d}^2 y}{\mathrm{d}x^2}-\frac{\mathrm{d}y}{\mathrm{d}x}-6y=50\sin x
\] \\
subject to the conditions $y(0)=3$ and $y'(0)=-6$. \marks{10}

\vspace*{10pt}

\begin{tcolorbox}[colback=white, colframe=scarlet, title={\textbf{Solution}}, breakable, grow to left by=6mm, grow to right by=10mm]
We solve
\[
\frac{\mathrm{d}^2y}{\mathrm{d}x^2}-\frac{\mathrm{d}y}{\mathrm{d}x}-6y=50\sin x
\]
subject to
\[
y(0)=3, \qquad y'(0)=-6
\]

First, solve the homogeneous equation
\[
\frac{\mathrm{d}^2y}{\mathrm{d}x^2}-\frac{\mathrm{d}y}{\mathrm{d}x}-6y=0
\]

Its auxiliary equation is
\[
\lambda^2-\lambda-6=0
\]

Factorising,
\[
(\lambda-3)(\lambda+2)=0
\]

So the roots are \(\lambda=3\) and \(\lambda=-2\). Hence the complementary function is
\[
y_c=Ae^{3x}+Be^{-2x}
\]

Now find a particular integral. Since the right-hand side is \(50\sin x\), try
\[
y_p=a\sin x+b\cos x
\]

Then
\begin{align*}
\frac{\mathrm{d}y_p}{\mathrm{d}x} &= a\cos x-b\sin x \\
\frac{\mathrm{d}^2y_p}{\mathrm{d}x^2} &= -a\sin x-b\cos x
\end{align*}

Substituting into the differential equation,
\begin{align*}
\frac{\mathrm{d}^2y_p}{\mathrm{d}x^2}-\frac{\mathrm{d}y_p}{\mathrm{d}x}-6y_p
&=\left(-a\sin x-b\cos x\right)-\left(a\cos x-b\sin x\right)-6\left(a\sin x+b\cos x\right) \\
&=-a\sin x-b\cos x-a\cos x+b\sin x-6a\sin x-6b\cos x \\
&=(b-7a)\sin x+(-a-7b)\cos x
\end{align*}

We want this to equal \(50\sin x\), so equating coefficients gives
\begin{align*}
b-7a&=50 \\
-a-7b&=0
\end{align*}

From the second equation,
\[
a=-7b
\]

Substitute into the first:
\begin{align*}
b-7(-7b)&=50 \\
b+49b&=50 \\
50b&=50 \\
b&=1
\end{align*}

So
\[
a=-7
\]

Therefore the particular integral is
\[
y_p=-7\sin x+\cos x
\]

Hence the general solution is
\[
y=Ae^{3x}+Be^{-2x}-7\sin x+\cos x
\]

Now use the initial conditions. First,
\begin{align*}
y(0)&=Ae^0+Be^0-7\sin 0+\cos 0 \\
&=A+B+1
\end{align*}

Since \(y(0)=3\),
\[
A+B+1=3
\]
so
\[
A+B=2
\]

Next differentiate the general solution:
\[
y'(x)=3Ae^{3x}-2Be^{-2x}-7\cos x-\sin x
\]

Now apply \(y'(0)=-6\):
\begin{align*}
y'(0)&=3Ae^0-2Be^0-7\cos 0-\sin 0 \\
&=3A-2B-7
\end{align*}

So
\[
3A-2B-7=-6
\]
hence
\[
3A-2B=1
\]

We now solve
\begin{align*}
A+B&=2 \\
3A-2B&=1
\end{align*}

Multiply the first equation by \(2\):
\[
2A+2B=4
\]

Add this to \(3A-2B=1\):
\begin{align*}
(2A+2B)+(3A-2B)&=4+1 \\
5A&=5 \\
A&=1
\end{align*}

Then from \(A+B=2\),
\begin{align*}
1+B&=2 \\
B&=1
\end{align*}

Therefore the required solution is
\[
y=e^{3x}+e^{-2x}-7\sin x+\cos x
\]
\end{tcolorbox}


\newpage

% ---- Question 8 ----
\questionitem
% [Source: _QBT___Solns__2nd_Order_Differential_Equations, Q8]

Determine the solution $y(x)$ of the initial value problem \\
\[
\frac{\mathrm{d}^2 y}{\mathrm{d}x^2} - 4\frac{\mathrm{d}y}{\mathrm{d}x} + 4y = xe^{2x}
\] \\
subject to
\[
 y(0)=1, \qquad \left.\frac{\mathrm{d}y}{\mathrm{d}x}\right|_{x=0} = 0
\] \marks[-20pt]{10}

\vspace*{10pt}

\begin{tcolorbox}[colback=white, colframe=scarlet, title={\textbf{Solution}}, breakable, grow to left by=6mm, grow to right by=10mm]
For the homogeneous equation,
\[
\frac{\mathrm{d}^2 y}{\mathrm{d}x^2}-4\frac{\mathrm{d}y}{\mathrm{d}x}+4y=0
\]

The auxiliary equation is
\[
\lambda^2-4\lambda+4=0
\]

So
\[
(\lambda-2)^2=0
\]

Hence the complementary function is
\[
y_c=(A+Bx)e^{2x}
\]

For a particular integral, since the right-hand side is $xe^{2x}$ and $\lambda=2$ is a repeated root, try
\[
y_p=e^{2x}(ax^3+bx^2)
\]

Then
\begin{align*}
\frac{\mathrm{d}y_p}{\mathrm{d}x}
&= \frac{\mathrm{d}}{\mathrm{d}x}\left[e^{2x}(ax^3+bx^2)\right] \\
&= 2e^{2x}(ax^3+bx^2)+e^{2x}(3ax^2+2bx) \\
&= e^{2x}\left(2ax^3+(2b+3a)x^2+2bx\right)
\end{align*}

Differentiate again:
\begin{align*}
\frac{\mathrm{d}^2y_p}{\mathrm{d}x^2}
&= \frac{\mathrm{d}}{\mathrm{d}x}\left[e^{2x}\left(2ax^3+(2b+3a)x^2+2bx\right)\right] \\
&= 2e^{2x}\left(2ax^3+(2b+3a)x^2+2bx\right) \\
&\qquad +e^{2x}\left(6ax^2+2(2b+3a)x+2b\right) \\
&= e^{2x}\left(4ax^3+(4b+12a)x^2+(8b+6a)x+2b\right)
\end{align*}

Substituting into the differential equation gives
\begin{align*}
\frac{\mathrm{d}^2y_p}{\mathrm{d}x^2}-4\frac{\mathrm{d}y_p}{\mathrm{d}x}+4y_p
&= e^{2x}\left(4ax^3+(4b+12a)x^2+(8b+6a)x+2b\right) \\
&\qquad -4e^{2x}\left(2ax^3+(2b+3a)x^2+2bx\right) \\
&\qquad +4e^{2x}(ax^3+bx^2)
\end{align*}

So
\[
\frac{\mathrm{d}^2y_p}{\mathrm{d}x^2}-4\frac{\mathrm{d}y_p}{\mathrm{d}x}+4y_p
=e^{2x}(6ax+2b)
\]

Comparing with $xe^{2x}$, we get
\[
6a=1, \qquad 2b=0
\]

So
\[
a=\frac{1}{6}, \qquad b=0
\]

Therefore
\[
y_p=\frac{x^3}{6}e^{2x}
\]

Thus the general solution is
\[
y=(A+Bx)e^{2x}+\frac{x^3}{6}e^{2x}
\]

So
\[
y=e^{2x}\left(A+Bx+\frac{x^3}{6}\right)
\]

Now use the initial conditions

From $y(0)=1$,
\begin{align*}
1
&= e^{0}\left(A+B\cdot 0+\frac{0^3}{6}\right) \\
&= A
\end{align*}

So
\[
A=1
\]

Now differentiate
\begin{align*}
\frac{\mathrm{d}y}{\mathrm{d}x}
&= \frac{\mathrm{d}}{\mathrm{d}x}\left[e^{2x}\left(A+Bx+\frac{x^3}{6}\right)\right] \\
&= 2e^{2x}\left(A+Bx+\frac{x^3}{6}\right)+e^{2x}\left(B+\frac{x^2}{2}\right) \\
&= e^{2x}\left(2A+2Bx+\frac{x^3}{3}+B+\frac{x^2}{2}\right)
\end{align*}

At $x=0$,
\[
\frac{\mathrm{d}y}{\mathrm{d}x}(0)=2A+B
\]

Since $\frac{\mathrm{d}y}{\mathrm{d}x}(0)=0$ and $A=1$,
\[
0=2+B
\]

Hence
\[
B=-2
\]

Substituting $A=1$ and $B=-2$ gives
\[
y(x)=e^{2x}\left(\frac{x^3}{6}-2x+1\right)
\]

Therefore, the solution of the initial value problem is
\[
y(x)=e^{2x}\left(\frac{x^3}{6}-2x+1\right)
\]
\end{tcolorbox}


\newpage

% ---- Question 9 ----
\questionitem
% [Source: _QBT___Solns__2nd_Order_Differential_Equations, Q9]

\begin{enumerate}[label=\textbf{(\alph*)}]
  \item Find the general solution of the differential equation
  \[
  \frac{\mathrm{d}^2 y}{\mathrm{d}x^2}+4y=0
  \]
  \marks[-20pt]{2}
  \item Hence, by considering a particular integral of the form $x(A\cos 2x+B\sin 2x)$, find the general solution of the differential equation
  \[
  \frac{\mathrm{d}^2 y}{\mathrm{d}x^2}+4y=8\cos 2x-12\sin 2x
  \]
  \marks[-20pt]{4}
\end{enumerate}

\vspace*{10pt}

\begin{tcolorbox}[colback=white, colframe=scarlet, title={\textbf{Solution}}, breakable, grow to left by=6mm, grow to right by=10mm]
\begin{enumerate}[label=\textbf{(\alph*)}]
  \item For
  \[
  \frac{\mathrm{d}^2 y}{\mathrm{d}x^2}+4y=0
  \]
  assume a solution of the form $y=e^{\lambda x}$. Then the auxiliary equation is
  \[
  \lambda^2+4=0
  \]
  So
  \[
  \lambda=\pm 2i
  \]
  Hence the general solution is
  \[
  y=C_1\cos 2x+C_2\sin 2x
  \]

  \item From part (a), the complementary solution is
  \[
  y_c=C_1\cos 2x+C_2\sin 2x
  \]
  Since the right-hand side involves $\cos 2x$ and $\sin 2x$, take a particular solution in the given form
  \[
  y_p=x(A\cos 2x+B\sin 2x)
  \]
  Differentiate:
  \begin{align*}
  \frac{\mathrm{d}y_p}{\mathrm{d}x}
  &= (A\cos 2x+B\sin 2x)+x(-2A\sin 2x+2B\cos 2x)
  \end{align*}
  Differentiate again:
  \begin{align*}
  \frac{\mathrm{d}^2 y_p}{\mathrm{d}x^2}
  &= (-2A\sin 2x+2B\cos 2x)+(-2A\sin 2x+2B\cos 2x) \\
  &\qquad {} + x(-4A\cos 2x-4B\sin 2x) \\
  &= -4A\sin 2x+4B\cos 2x-4x(A\cos 2x+B\sin 2x)
  \end{align*}
  Now substitute into $\dfrac{\mathrm{d}^2 y}{\mathrm{d}x^2}+4y$:
  \begin{align*}
  \frac{\mathrm{d}^2 y_p}{\mathrm{d}x^2}+4y_p
  &= \left[-4A\sin 2x+4B\cos 2x-4x(A\cos 2x+B\sin 2x)\right] \\
  &\qquad {} + 4x(A\cos 2x+B\sin 2x) \\
  &= 4B\cos 2x-4A\sin 2x
  \end{align*}
  This must equal the right-hand side:
  \[
  4B\cos 2x-4A\sin 2x=8\cos 2x-12\sin 2x
  \]
  Equating coefficients gives
  \begin{align*}
  4B &= 8 \\
  -4A &= -12
  \end{align*}
  so
  \[
  B=2,\qquad A=3
  \]
  Therefore
  \[
  y_p=x(3\cos 2x+2\sin 2x)
  \]
  Hence the general solution is
  \[
  y=y_c+y_p=C_1\cos 2x+C_2\sin 2x+x(3\cos 2x+2\sin 2x)
  \]
\end{enumerate}
\end{tcolorbox}


\newpage

% ---- Question 10 ----
\questionitem
% [Source: _QBT___Solns__2nd_Order_Differential_Equations, Q10]

Find the particular solution of the differential equation\\

\[
\frac{\mathrm{d}^2 y}{\mathrm{d}t^2}+4\frac{\mathrm{d}y}{\mathrm{d}t}+13y=28\cos 3t-4\sin 3t
\]\\
given that, when $t=0$, $y=5$ and $\frac{\mathrm{d}y}{\mathrm{d}t}=7 \marks{10}$ \\

\vspace*{10pt}

\begin{tcolorbox}[colback=white, colframe=scarlet, title={\textbf{Solution}}, breakable, grow to left by=6mm, grow to right by=10mm]
We solve
\[
\frac{\mathrm{d}^2y}{\mathrm{d}t^2}+4\frac{\mathrm{d}y}{\mathrm{d}t}+13y=28\cos 3t-4\sin 3t
\]
by finding the complementary function and a particular integral, then using the initial conditions.

For the complementary function, solve the auxiliary equation
\[
\lambda^2+4\lambda+13=0
\]
Using the quadratic formula,
\begin{align*}
\lambda&=\frac{-4\pm\sqrt{4^2-4(1)(13)}}{2} \\[4pt]
&=-2\pm 3\mathrm{i}
\end{align*}
Hence the complementary function is
\[
y_c=\mathrm{e}^{-2t}(A\cos 3t+B\sin 3t)
\]

For the particular integral, since the right-hand side is of the form $28\cos 3t-4\sin 3t$, try
\[
y_p=a\cos 3t+b\sin 3t
\]
Then
\begin{align*}
\frac{\mathrm{d}y_p}{\mathrm{d}t}&=-3a\sin 3t+3b\cos 3t \\
\frac{\mathrm{d}^2y_p}{\mathrm{d}t^2}&=-9a\cos 3t-9b\sin 3t
\end{align*}
Substituting into the differential equation,
\begin{align*}
\frac{\mathrm{d}^2y_p}{\mathrm{d}t^2}+4\frac{\mathrm{d}y_p}{\mathrm{d}t}+13y_p
&=\left(-9a\cos 3t-9b\sin 3t\right)+4\left(-3a\sin 3t+3b\cos 3t\right)+13\left(a\cos 3t+b\sin 3t\right) \\
&=\left(-9a+12b+13a\right)\cos 3t+\left(-9b-12a+13b\right)\sin 3t \\
&=(4a+12b)\cos 3t+(4b-12a)\sin 3t
\end{align*}
Equating coefficients with $28\cos 3t-4\sin 3t$ gives
\begin{align*}
4a+12b&=28 \\
4b-12a&=-4
\end{align*}
Dividing both equations by $4$,
\begin{align*}
a+3b&=7 \\
b-3a&=-1
\end{align*}
From the second equation,
\[
b=3a-1
\]
Substitute into the first:
\begin{align*}
a+3(3a-1)&=7 \\
a+9a-3&=7 \\
10a&=10 \\
a&=1
\end{align*}
Then
\[
b=3(1)-1=2
\]
So
\[
y_p=\cos 3t+2\sin 3t
\]

Hence the general solution is
\[
y=\mathrm{e}^{-2t}(A\cos 3t+B\sin 3t)+\cos 3t+2\sin 3t
\]

Now use the initial conditions.

First, when $t=0$, $y=5$:
\begin{align*}
5&=\mathrm{e}^{0}(A\cos 0+B\sin 0)+\cos 0+2\sin 0 \\
&=A(1)+B(0)+1+0 \\
&=A+1
\end{align*}
So
\[
A=4
\]

Next differentiate $y$.

For the homogeneous part,
\begin{align*}
\frac{\mathrm{d}}{\mathrm{d}t}\left[\mathrm{e}^{-2t}(A\cos 3t+B\sin 3t)\right]
&=\mathrm{e}^{-2t}\left(-2\right)(A\cos 3t+B\sin 3t)+\mathrm{e}^{-2t}\left(-3A\sin 3t+3B\cos 3t\right) \\
&=\mathrm{e}^{-2t}\left[(3B-2A)\cos 3t+(-3A-2B)\sin 3t\right]
\end{align*}
Also,
\[
\frac{\mathrm{d}}{\mathrm{d}t}\left(\cos 3t+2\sin 3t\right)=-3\sin 3t+6\cos 3t
\]
Therefore
\[
\frac{\mathrm{d}y}{\mathrm{d}t}=\mathrm{e}^{-2t}\left[(3B-2A)\cos 3t+(-3A-2B)\sin 3t\right]-3\sin 3t+6\cos 3t
\]

Using $\frac{\mathrm{d}y}{\mathrm{d}t}=7$ when $t=0$,
\begin{align*}
7&=\mathrm{e}^{0}\left[(3B-2A)\cos 0+(-3A-2B)\sin 0\right]-3\sin 0+6\cos 0 \\
&=(3B-2A)(1)+0+0+6 \\
&=3B-2A+6
\end{align*}
Substitute $A=4$:
\begin{align*}
7&=3B-8+6 \\
7&=3B-2 \\
3B&=9 \\
B&=3
\end{align*}

Hence the required particular solution is
\[
y=(4\cos 3t+3\sin 3t)\mathrm{e}^{-2t}+\cos 3t+2\sin 3t
\]
\end{tcolorbox}


\newpage

% ---- Question 11 ----
\questionitem
% [Source: _QBT___Solns__2nd_Order_Differential_Equations, Q11]

Find the particular solution of the differential equation\\
\[
\frac{\mathrm{d}^2 y}{\mathrm{d}t^2}+4\frac{\mathrm{d}y}{\mathrm{d}t}+5y=8\cos t
\]\\
given that, when $t=0$, $y=3$ and $\frac{\mathrm{d}y}{\mathrm{d}t}=1 \marks{10}$ \\

\vspace*{10pt}

\begin{tcolorbox}[colback=white, colframe=scarlet, title={\textbf{Solution}}, breakable, grow to left by=6mm, grow to right by=10mm]
First solve the associated homogeneous equation
\[
\frac{\mathrm{d}^2y}{\mathrm{d}t^2}+4\frac{\mathrm{d}y}{\mathrm{d}t}+5y=0
\]

Its auxiliary equation is
\[
\lambda^2+4\lambda+5=0
\]

Using the quadratic formula:
\begin{align*}
\lambda&=\frac{-4\pm\sqrt{4^2-4(1)(5)}}{2(1)}\\[4pt]
&=-2\pm \mathrm{i}
\end{align*}

Hence the complementary function is
\[
y_c=e^{-2t}(C_1\cos t+C_2\sin t)
\]

Now find a particular integral for
\[
\frac{\mathrm{d}^2y}{\mathrm{d}t^2}+4\frac{\mathrm{d}y}{\mathrm{d}t}+5y=8\cos t
\]

Since the forcing term is \(8\cos t\), try
\[
y_p=A\cos t+B\sin t
\]

Then
\[
\frac{\mathrm{d}y_p}{\mathrm{d}t}=-A\sin t+B\cos t
\]
and
\[
\frac{\mathrm{d}^2y_p}{\mathrm{d}t^2}=-A\cos t-B\sin t
\]

Substituting into the differential equation:
\begin{align*}
\frac{\mathrm{d}^2y_p}{\mathrm{d}t^2}+4\frac{\mathrm{d}y_p}{\mathrm{d}t}+5y_p
&=(-A\cos t-B\sin t)+4(-A\sin t+B\cos t)+5(A\cos t+B\sin t)\\
&=(-A+4B+5A)\cos t+(-B-4A+5B)\sin t\\
&=(4A+4B)\cos t+(-4A+4B)\sin t
\end{align*}

This must equal \(8\cos t\), so equating coefficients gives
\[
4A+4B=8,\qquad -4A+4B=0
\]

Divide both equations by \(4\):
\[
A+B=2,\qquad -A+B=0
\]

From \(-A+B=0\), we get
\[
B=A
\]

Substitute into \(A+B=2\):
\[
A+A=2
\]
\[
2A=2
\]
\[
A=1
\]

Therefore
\[
B=1
\]

So the particular integral is
\[
y_p=\cos t+\sin t
\]

Hence the general solution is
\[
y=e^{-2t}(C_1\cos t+C_2\sin t)+\cos t+\sin t
\]

Now use the initial conditions.

Given \(y(0)=3\), substitute \(t=0\):
\begin{align*}
3&=e^{0}(C_1\cos 0+C_2\sin 0)+\cos 0+\sin 0\\
&=1(C_1\cdot 1+C_2\cdot 0)+1+0\\
&=C_1+1
\end{align*}

So
\[
C_1+1=3
\]
\[
C_1=2
\]

Next differentiate the general solution:
\begin{align*}
\frac{\mathrm{d}y}{\mathrm{d}t}
&=\frac{\mathrm{d}}{\mathrm{d}t}\left[e^{-2t}(C_1\cos t+C_2\sin t)\right]+\frac{\mathrm{d}}{\mathrm{d}t}(\cos t+\sin t)\\
&=-2e^{-2t}(C_1\cos t+C_2\sin t)+e^{-2t}(-C_1\sin t+C_2\cos t)-\sin t+\cos t\\
&=e^{-2t}\left[(-2C_1+C_2)\cos t+(-C_1-2C_2)\sin t\right]-\sin t+\cos t
\end{align*}

Given \(\frac{\mathrm{d}y}{\mathrm{d}t}(0)=1\), substitute \(t=0\):
\begin{align*}
1&=e^{0}\left[(-2C_1+C_2)\cos 0+(-C_1-2C_2)\sin 0\right]-\sin 0+\cos 0\\
&=1\left[(-2C_1+C_2)\cdot 1+(-C_1-2C_2)\cdot 0\right]-0+1\\
&=(-2C_1+C_2)+1
\end{align*}

So
\[
(-2C_1+C_2)+1=1
\]
\[
-2C_1+C_2=0
\]

Using \(C_1=2\):
\[
-2(2)+C_2=0
\]
\[
-4+C_2=0
\]
\[
C_2=4
\]

Therefore the particular solution satisfying the given conditions is
\[
y=e^{-2t}(2\cos t+4\sin t)+\cos t+\sin t
\]
\end{tcolorbox}


\newpage

% ---- Question 12 ----
\questionitem
% [Source: _QBT___Solns__2nd_Order_Differential_Equations, Q12]

\begin{enumerate}[label=\textbf{(\alph*)}]
  \item Show that, for all real $x$, \\
  \[
  (1+\tanh x)\cosh x = e^{x}
  \]
  \marks[-20pt]{2} \\
  \item Hence solve the differential equation \\
  \[
  \frac{\mathrm{d}^2 y}{\mathrm{d}x^2} - 2\frac{\mathrm{d}y}{\mathrm{d}x} + y = 3(1+\tanh x)\cosh x
  \]\\
  given that, when $x=0$, $y=1$ and $\frac{\mathrm{d}y}{\mathrm{d}x}=0$ \marks{10} \\
\end{enumerate}

\vspace*{10pt}

\begin{tcolorbox}[colback=white, colframe=scarlet, title={\textbf{Solution}}, breakable, grow to left by=6mm, grow to right by=10mm]
\begin{enumerate}[label=\textbf{(\alph*)}]
  \item Using \(\tanh x=\dfrac{\sinh x}{\cosh x}\),
\begin{align*}
(1+\tanh x)\cosh x
&=\left(1+\frac{\sinh x}{\cosh x}\right)\cosh x \\
&=\cosh x+\sinh x \\
&=\frac{e^x+e^{-x}}{2}+\frac{e^x-e^{-x}}{2} \\
&=e^x
\end{align*}
Hence, for all real \(x\),
\[
(1+\tanh x)\cosh x=e^x
\]

\item From part \((a)\), the differential equation becomes
\[
\frac{\mathrm{d}^2 y}{\mathrm{d}x^2}-2\frac{\mathrm{d}y}{\mathrm{d}x}+y=3e^x
\]

For the homogeneous equation,
\[
\frac{\mathrm{d}^2 y}{\mathrm{d}x^2}-2\frac{\mathrm{d}y}{\mathrm{d}x}+y=0
\]

The auxiliary equation is
\[
\lambda^2-2\lambda+1=0
\]

So
\[
(\lambda-1)^2=0
\]

Hence the complementary function is
\[
y_c=(A+Bx)e^x
\]

For a particular integral, since the right-hand side is $3e^x$ and $\lambda=1$ is a repeated root, try
\[
y_p=Cx^2e^x
\]

Then
\begin{align*}
\frac{\mathrm{d}y_p}{\mathrm{d}x}
&= \frac{\mathrm{d}}{\mathrm{d}x}\left(Cx^2e^x\right) \\
&= 2Cxe^x+Cx^2e^x \\
&= e^x\left(Cx^2+2Cx\right)
\end{align*}

Differentiate again:
\begin{align*}
\frac{\mathrm{d}^2y_p}{\mathrm{d}x^2}
&= \frac{\mathrm{d}}{\mathrm{d}x}\left[e^x\left(Cx^2+2Cx\right)\right] \\
&= e^x\left(Cx^2+2Cx\right)+e^x\left(2Cx+2C\right) \\
&= e^x\left(Cx^2+4Cx+2C\right)
\end{align*}

Substituting into the differential equation gives
\begin{align*}
\frac{\mathrm{d}^2y_p}{\mathrm{d}x^2}-2\frac{\mathrm{d}y_p}{\mathrm{d}x}+y_p
&= e^x\left(Cx^2+4Cx+2C\right)-2e^x\left(Cx^2+2Cx\right)+Cx^2e^x \\
&= 2Ce^x
\end{align*}

So
\[
2Ce^x=3e^x
\]

Hence
\[
2C=3
\]

So
\[
C=\frac{3}{2}
\]

Therefore
\[
y_p=\frac{3x^2}{2}e^x
\]

Thus the general solution is
\[
y=(A+Bx)e^x+\frac{3x^2}{2}e^x
\]

So
\[
y=e^x\left(A+Bx+\frac{3x^2}{2}\right)
\]

Now use the initial conditions

From $y(0)=1$,
\begin{align*}
1
&= e^0\left(A+B\cdot 0+\frac{3(0)^2}{2}\right) \\
&= A
\end{align*}

So
\[
A=1
\]

Now differentiate
\begin{align*}
\frac{\mathrm{d}y}{\mathrm{d}x}
&= \frac{\mathrm{d}}{\mathrm{d}x}\left[e^x\left(A+Bx+\frac{3x^2}{2}\right)\right] \\
&= e^x\left(A+Bx+\frac{3x^2}{2}\right)+e^x(B+3x) \\
&= e^x\left(A+Bx+\frac{3x^2}{2}+B+3x\right)
\end{align*}

At $x=0$,
\[
\frac{\mathrm{d}y}{\mathrm{d}x}(0)=A+B
\]

Since $\frac{\mathrm{d}y}{\mathrm{d}x}(0)=0$ and $A=1$,
\[
0=1+B
\]

Hence
\[
B=-1
\]

Substituting $A=1$ and $B=-1$ into the expression for $y$,
\[
y=e^x\left(\frac{3x^2}{2}-x+1\right)
\]

Therefore the required solution is
\[
y=e^x\left(\frac{3x^2}{2}-x+1\right)
\]

\end{enumerate}
\end{tcolorbox}
\end{enumerate}
\end{document}
